% ============================================================
% Random Partial Benders Decomposition for Robust CFLP
% Adapted from Pauphilet & Wang (2024) to the two-stage
% robust Congested Facility Location Problem with M/M/1
% congestion and budget uncertainty.
% ============================================================

\section{Scalable C\&CG via Partial Scenario Blocks}
\label{sec:rpbd}

\subsection{Motivation}
\label{sec:rpbd:motivation}

The standard Column-and-Constraint Generation (C\&CG) algorithm
of \citet{ZengZhao2013} solves the two-stage robust problem
\begin{equation}
  \min_{x \in \mathcal{X}}\; f(x) + \max_{\varepsilon \in \Xi}\; Q(x,\varepsilon)
  \label{eq:robust}
\end{equation}
by iteratively adding \emph{scenario blocks} to a master Mixed-Integer
Second-Order Cone Programme (MISOCP).
Each scenario block introduced at iteration $t$ for scenario
$\bar{\varepsilon}^{t}$ adds
\begin{equation*}
  O\!\left(|\mathcal{I}|\cdot|\mathcal{J}|\cdot|\mathcal{R}|\right)
\end{equation*}
continuous variables and constraints to the master, where
$\mathcal{I}$, $\mathcal{J}$, $\mathcal{R}$ are the sets of
customers, candidate facilities, and capacity levels, respectively.
After $T$ iterations the master contains
$O\!\left(T \cdot |\mathcal{I}|\cdot|\mathcal{J}|\cdot|\mathcal{R}|\right)$
variables, making each subsequent solve progressively more expensive.
Table~\ref{tab:time performance} confirms this: for instances with
$|\mathcal{J}| \ge 15$ the algorithm either times out or requires
hours of computation even for moderate $\Gamma$.

To address this scalability bottleneck we propose a
\emph{partial block} strategy inspired by the Random Partial Benders
Decomposition (RPBD) framework of \citet{PauphiletWang2024}.
The key idea is to add, at each iteration, \emph{two} complementary
but cheap cuts instead of one expensive full block:
\begin{enumerate}
  \item a \textbf{linear Benders cut} covering \emph{all} active
        facilities — computationally cheap, globally valid, but loose;
  \item a \textbf{partial SOC block} covering only the
        $k = \lceil \phi\,|\mathcal{J}^{t}| \rceil$ most
        \emph{congestion-critical} facilities — tighter for the
        facilities that matter most, at a fraction of the cost.
\end{enumerate}
Together these two cuts enforce
$\eta \ge \max\!\left(
  \text{Benders}(x),\;
  \text{SOC}_{\mathcal{J}^{t}_{\mathrm{retain}}}(x,y,C)
\right)$
without an auxiliary binary variable, because both are added as
separate $\ge$ constraints on the surrogate objective variable $\eta$.
The master grows at $O(\phi \cdot |\mathcal{I}|\cdot|\mathcal{J}|\cdot|\mathcal{R}|)$
per iteration — a $1/\phi$ reduction in variable growth.

% ------------------------------------------------------------
\subsection{Background: Standard C\&CG Master}
\label{sec:rpbd:standard}

Let $x \in \{0,1\}^{|\mathcal{J}|\times|\mathcal{R}|}$ denote the
first-stage facility-opening decisions and $\eta \in \mathbb{R}$ the
epigraph variable bounding worst-case recourse cost.
After $T$ iterations the standard C\&CG master is
\begin{align}
  \min_{x,\eta,y^t,C^t,Q^t}\;&
    \sum_{j\in\mathcal{J}}\sum_{r\in\mathcal{R}} f_{jr}\,x_{jr}
    + \eta
    \label{eq:master:obj}\\
  \text{s.t.}\quad&
    \sum_{r\in\mathcal{R}} x_{jr} \le 1
    && \forall\, j\in\mathcal{J}
    \label{eq:master:open}\\
  &\eta \ge
    \sum_{i,j,r} c^1_{ij}\,y^t_{ijr}
    + \sum_i c^2_i\,Q^t_i
    + \sum_{j,r} w_j\,C^t_{jr}
    && \forall\, t = 1,\ldots,T
    \label{eq:master:optcut}\\
  &\text{(assignment, capacity, M/M/1 SOC constraints for each block } t)
    \notag\\
  &x_{jr}\in\{0,1\},\quad
    y^t_{ijr}\in[0,1],\quad
    Q^t_i,C^t_{jr}\ge 0
    \notag
\end{align}
where $c^1_{ij}$, $c^2_i$, $w_j$ are transportation, congestion-penalty,
and congestion-cost coefficients respectively.
The M/M/1 second-order cone constraints for scenario $t$ are
\begin{equation}
  \left(\lambda^t_{jr}\right)^2
  \;\le\;
  \Bigl(\bar\varepsilon^t_j\,\mu_{jr}\,C^t_{jr} - \lambda^t_{jr}\Bigr)
  \Bigl(\bar\varepsilon^t_j\,\mu_{jr}\,x_{jr}   - \lambda^t_{jr}\Bigr),
  \qquad \forall\,j\in\mathcal{J}^t,\; r\in\mathcal{R}
  \label{eq:master:soc}
\end{equation}
where $\lambda^t_{jr} = \sum_i d_i\,y^t_{ijr}$ is the load at facility
$j$ under capacity level $r$, $\mu_{jr}$ is the capacity, and
$\bar\varepsilon^t_j = \sum_h (1 - h/(H-1))\,\bar\varepsilon^t_{jh}$
is the effective capacity factor under scenario $\bar\varepsilon^t$.

% ------------------------------------------------------------
\subsection{Benders Cut from the Subproblem Dual}
\label{sec:rpbd:benders}

The subproblem separation oracle maximises over $\varepsilon\in\Xi$
and simultaneously dualises the inner SOCP recourse.
When solved with dual extraction (\texttt{return\_duals=True}),
it returns dual variables
$\alpha_i,\,\beta_i,\,t2_i$ for customer constraints and
linearised products
$\tilde\theta_{jh} = \varepsilon_{jh}\,\theta_j$,
$\tilde u_{2,jh} = \varepsilon_{jh}\,u_{2,j}$,
$\tilde\gamma_{jh} = \varepsilon_{jh}\,\gamma_j$
for facility-level cone multipliers.
The resulting Benders optimality cut is
\begin{equation}
  \eta \;\ge\;
  -\sum_{i\in\mathcal{I}} \alpha_i
  \;-\; \sum_{j\in\mathcal{J}}\sum_{r\in\mathcal{R}}\sum_{h=0}^{H-1}
        \mu_{jr}\,x_{jr}\,\Bigl(1 - \tfrac{h}{H-1}\Bigr)
        \Bigl(\tilde\theta_{jh} + \tilde u_{2,jh} + \tilde\gamma_{jh}\Bigr)
  \;+\; \sum_{i\in\mathcal{I}} t2_i
  \;-\; \sum_{i\in\mathcal{I}} \beta_i.
  \label{eq:benders:cut}
\end{equation}
This cut is \emph{linear} in $x$ and provides a valid global lower
bound on $Q(x,\bar\varepsilon)$ for \emph{all} $j\in\mathcal{J}$.
Its weakness is that it approximates the nonlinear M/M/1 congestion
cost with a supporting hyperplane, which can be loose when the
current $x$ is far from the dual solution.

% ------------------------------------------------------------
\subsection{Partial SOC Block Construction}
\label{sec:rpbd:partial}

Given scenario $\bar\varepsilon^t$ with active set
$\mathcal{J}^t = \{j : \bar\varepsilon^t_j > 0\}$,
define the \emph{criticality score} of facility $j$ as
\begin{equation}
  \sigma_j^t
  \;=\;
  \frac{w_j \cdot \bar\varepsilon^t_j}
       {\bar\mu_j + \delta},
  \qquad
  \bar\mu_j = \frac{1}{|\mathcal{R}|}\sum_{r\in\mathcal{R}} \mu_{jr},
  \label{eq:score}
\end{equation}
where $w_j$ is the congestion cost, $\bar\varepsilon^t_j$ is the
effective capacity degradation, and $\delta > 0$ is a small constant
for numerical stability.
Facilities with high congestion cost and high degradation under
$\bar\varepsilon^t$ contribute most to the worst-case recourse
and are therefore prioritised for exact SOC treatment.

Let $\phi \in (0,1)$ be the \emph{retention fraction} (parameter
\texttt{partial\_block\_fraction}).
Define the \emph{retained set}
\begin{equation}
  \mathcal{J}^t_{\mathrm{retain}}
  \;=\;
  \text{top-}k\text{ facilities in } \mathcal{J}^t
  \text{ by score } \sigma_j^t,
  \quad
  k = \left\lceil \phi\,|\mathcal{J}^t| \right\rceil.
  \label{eq:retain}
\end{equation}
The partial SOC block for $\mathcal{J}^t_{\mathrm{retain}}$ adds
variables $y^t_{ijr}$, $C^t_{jr}$, $V^t_{1,jr}$, $V^t_{2,jr}$,
$V^t_{3,jr}$, $Q^t_i$ and the constraints
\eqref{eq:master:optcut}--\eqref{eq:master:soc}
restricted to $j \in \mathcal{J}^t_{\mathrm{retain}}$.
Its associated opt-cut is
\begin{equation}
  \eta \;\ge\;
  \sum_{i\in\mathcal{I}}\sum_{j\in\mathcal{J}^t_{\mathrm{retain}}}\sum_{r\in\mathcal{R}}
    c^1_{ij}\,y^t_{ijr}
  + \sum_{i\in\mathcal{I}} c^2_i\,Q^t_i
  + \sum_{j\in\mathcal{J}^t_{\mathrm{retain}}}\sum_{r\in\mathcal{R}} w_j\,C^t_{jr}.
  \label{eq:partial:optcut}
\end{equation}
Since all costs are non-negative,
$Q(x,\bar\varepsilon^t) \ge Q_{\mathcal{J}^t_{\mathrm{retain}}}(x,\bar\varepsilon^t)$,
so \eqref{eq:partial:optcut} is a valid lower bound on the true
worst-case recourse even though it covers only a subset of facilities.

\begin{proposition}[Validity of the partial block]
\label{prop:validity}
Both the Benders cut \eqref{eq:benders:cut} and the partial SOC
opt-cut \eqref{eq:partial:optcut} are valid lower bounds on
$Q(x,\bar\varepsilon^t)$ for all feasible $x \in \mathcal{X}$.
Adding both as separate $\ge$ constraints on $\eta$ enforces
$\eta \ge \max\!\left(\eqref{eq:benders:cut},\,\eqref{eq:partial:optcut}\right)$,
which is a tighter bound than either cut alone.
\end{proposition}

\begin{proof}
The Benders cut follows from LP strong duality applied to the inner
SOCP for fixed $(x,\bar\varepsilon^t)$ \citep[Lemma~1]{ZengZhao2013}.
The partial opt-cut holds because
$Q(x,\bar\varepsilon^t) \ge \sum_{j\in\mathcal{J}^t_{\mathrm{retain}}} Q_j(x,\bar\varepsilon^t) \ge 0$
since individual facility costs are non-negative.
The $\max$ follows immediately from having both constraints active on
the same variable $\eta$. \qed
\end{proof}

% ------------------------------------------------------------
\subsection{Algorithm}
\label{sec:rpbd:algorithm}

Algorithm~\ref{alg:rpbd} states the complete procedure.

\begin{algorithm}[t]
\caption{Partial-Block C\&CG (RPBD mode)}
\label{alg:rpbd}
\begin{algorithmic}[1]
\Require
  Instance data, uncertainty budget $\Gamma$, levels $H$,
  initial solution $x^0$, tolerance $\epsilon_{\mathrm{tol}}$,
  retention fraction $\phi\in(0,1)$,
  master MIP gap $\delta_{\mathrm{mp}}$, master time limit $\tau$.
\Ensure
  Near-optimal robust solution $x^*$, bounds $\mathrm{UB}$, $\mathrm{LB}$.
\State $\mathrm{UB} \leftarrow +\infty$; \;
       $\mathrm{LB} \leftarrow -\infty$; \;
       $t \leftarrow 0$
\State Build master \eqref{eq:master:obj}--\eqref{eq:master:open}
       with variables $x$ and $\eta$ only.
\Repeat
  \State \textbf{// Subproblem (separation oracle with dual extraction)}
  \State $(\bar\varepsilon^t,\, \alpha, \beta, t_2, \tilde\theta, \tilde u_2, \tilde\gamma,\, Q^*) \leftarrow
         \textsc{Subproblem}(x^t,\,\text{return\_duals}=\mathrm{True})$
  \State $\mathrm{UB} \leftarrow \min\!\left(\mathrm{UB},\;
         Q^* + f(x^t)\right)$
  \State \textbf{// Add Benders cut \eqref{eq:benders:cut} (all $j\in\mathcal{J}$)}
  \State Add constraint \eqref{eq:benders:cut} to master.
  \State \textbf{// Add partial SOC block \eqref{eq:partial:optcut}--\eqref{eq:master:soc}}
  \State Compute scores $\sigma_j^t$ via \eqref{eq:score} for
         $j\in\mathcal{J}^t$.
  \State $\mathcal{J}^t_{\mathrm{retain}} \leftarrow
         \text{top-}\lceil\phi|\mathcal{J}^t|\rceil
         \text{ facilities by } \sigma_j^t$
  \State Add variables $y^t_{ijr}$, $C^t_{jr}$, $V^t_{\cdot,jr}$,
         $Q^t_i$ and constraints for $j\in\mathcal{J}^t_{\mathrm{retain}}$,
         including opt-cut \eqref{eq:partial:optcut} and SOC
         \eqref{eq:master:soc}.
  \State \textbf{// Solve master}
  \State $(x^{t+1},\,\eta^{t+1},\,\mathrm{LB}) \leftarrow
         \textsc{SolveMaster}(\delta_{\mathrm{mp}},\,\tau)$
  \State $t \leftarrow t + 1$
\Until{$(\mathrm{UB} - \mathrm{LB})\,/\,|\mathrm{UB}| \le \epsilon_{\mathrm{tol}}$
       \textbf{ or } time limit exceeded}
\State \Return $x^t$, $\mathrm{UB}$, $\mathrm{LB}$
\end{algorithmic}
\end{algorithm}

\begin{remark}[Master size growth]
At iteration $t$, the master contains
\begin{equation*}
  |\mathcal{J}||\mathcal{R}| + 1
  \;+\; t\!\cdot\!\left(
    |\mathcal{I}| + \lceil\phi|\mathcal{J}|\rceil \cdot |\mathcal{R}|
    \cdot(|\mathcal{I}|+4)
  \right)
\end{equation*}
variables, compared with
$|\mathcal{J}||\mathcal{R}| + 1 + t\cdot|\mathcal{I}|\cdot|\mathcal{J}|\cdot|\mathcal{R}|\cdot(1+4/|\mathcal{I}|)$
for standard C\&CG\@.
The reduction factor is approximately $\phi$ in the dominant term.
\end{remark}

\begin{remark}[Correctness]
Algorithm~\ref{alg:rpbd} is a valid outer-approximation scheme.
Each master solve yields a lower bound $\mathrm{LB} \le \min_x f(x)+\max_\varepsilon Q(x,\varepsilon)$
because both cuts \eqref{eq:benders:cut} and \eqref{eq:partial:optcut}
are valid relaxations.
The subproblem provides an upper bound $\mathrm{UB}$ by evaluating
the true worst-case recourse at the current $x^t$.
Convergence is guaranteed when $\phi > 0$ because the Benders cuts
alone are sufficient for asymptotic convergence (as in standard BDCP),
and the partial SOC blocks accelerate gap closure.
\end{remark}

% ------------------------------------------------------------
\subsection{Connection to RPBD}
\label{sec:rpbd:connection}

\citet{PauphiletWang2024} propose to randomly retain a fraction of
continuous recourse variables in the master of a Benders decomposition,
rather than projecting all of them out.
They show empirically that this typically halves convergence time on
network design and facility location benchmarks, and provide theoretical
bounds on the improvement to the LP relaxation.

Our adaptation differs in three respects suited to the CFLP setting:
\begin{enumerate}
  \item \textbf{Deterministic selection.}
        Rather than random retention, we select facilities
        greedily by criticality score $\sigma_j^t$
        \eqref{eq:score}. This is well-suited to our problem
        because the congestion bottleneck is concentrated at a
        small number of high-demand, low-capacity facilities.
  \item \textbf{Dual Benders cut always added.}
        Standard RPBD adds only the retained-variable block.
        We additionally add a full linear Benders cut at every
        iteration, providing a global lower bound even for the
        $1-\phi$ fraction of facilities not in the SOC block.
  \item \textbf{SOCP recourse.}
        The original RPBD paper targets LP recourse.
        Our recourse is a second-order cone programme due to M/M/1
        congestion costs; the partial-retention principle applies
        unchanged because both cuts remain valid lower bounds
        (Proposition~\ref{prop:validity}).
\end{enumerate}

% ------------------------------------------------------------
\subsection{Implementation Details}
\label{sec:rpbd:impl}

The algorithm is implemented in \texttt{rcflp/ccg.py} as
\texttt{solve\_CCG} with parameter
\texttt{partial\_block\_fraction} $= \phi \in (0,1)$.
When $\phi = 1$ the method reduces to standard C\&CG\@.
Key implementation choices:

\paragraph{Facility ranking.}
At each iteration, $\mathcal{J}^t_{\mathrm{retain}}$ is recomputed
from the current scenario $\bar\varepsilon^t$, so the selected
facilities adapt to which locations the adversary targets.

\paragraph{Dual extraction.}
Setting \texttt{return\_duals=True} in the subproblem call returns
the linearised dual products $\tilde\theta_{jh}$, $\tilde u_{2,jh}$,
$\tilde\gamma_{jh}$ needed for cut \eqref{eq:benders:cut}.
This is incompatible with the \texttt{exclude\_scenario} integer cut
used when \texttt{n\_scenarios=2}; accordingly \texttt{n\_scenarios}
is forced to 1 in RPBD mode.

\paragraph{Combination with block dropping.}
Setting \texttt{max\_active\_blocks} $= m$ additionally removes the
oldest non-binding partial block from the master once more than $m$
blocks are active, keeping master size bounded at
$O(m\cdot\phi\cdot|\mathcal{I}|\cdot|\mathcal{J}|\cdot|\mathcal{R}|)$
regardless of iteration count.

\paragraph{Recommended parameter values.}
Based on preliminary experiments we recommend $\phi = 0.4$ as a
starting point, retaining approximately the top 40\% of
congestion-critical facilities per block.
Values $\phi < 0.3$ risk slow convergence; values $\phi > 0.7$
provide little size reduction over the full block.

% ============================================================
% References (add to your main .bib file)
% ============================================================
%
% @article{ZengZhao2013,
%   author  = {Zeng, B. and Zhao, L.},
%   title   = {Solving two-stage robust optimization problems using
%              a column-and-constraint generation method},
%   journal = {Operations Research Letters},
%   volume  = {41},
%   number  = {5},
%   pages   = {457--461},
%   year    = {2013}
% }
%
% @article{PauphiletWang2024,
%   author  = {Pauphilet, J. and Wang, Y.},
%   title   = {The Surprising Performance of Random Partial Benders
%              Decomposition},
%   journal = {Preprint, London Business School},
%   year    = {2024}
% }
